% Proved unnumbered input for tower A.4.57.  Equation prefix CONE is unique.
\paragraph{A finite cone-fusion input with its characteristic source.}
Put
\[
 E=\langle a_1,a_2,a_3,u_1,u_2,u_3,u_4\mid
                         a_1^2=a_2^2=a_3^2=1\rangle
            =C_2*C_2*C_2*F_4.
                                                               \tag{CONE.1}
\]
Let $q:E\twoheadrightarrow D$ be a specified actual finite quotient
whose kernel lies in $E'E^2$; equivalently its actual quotient map
retains $H_1(E,\mathbb F_2)=\mathbb F_2^7$ with these named classes.
Define the ordinary subgroups and integers
\[
 \begin{gathered}
 N=\ker q,\qquad N_1=N'N^2,\qquad N_2=N_1'N_1^2,\qquad d=|D|,\\
 r_0=1+9d/2,\qquad r_1=1+2^{r_0}(r_0-1),\qquad
 m=d\,2^{r_0+r_1},\qquad A=E/\kappa_m(E).
 \end{gathered}                                                   \tag{CONE.2}
\]
The actual quotient maps are
\[
                  A\longrightarrow E/N_2\longrightarrow D.
                                                               \tag{CONE.3}
\]
Every order-two element of $E/N_2$ projects into a conjugate of one
of the three subgroups $q\langle a_j\rangle$ of $D$.
Moreover, for every actual automorphism $\alpha$ of the characteristic
finite source $A$, there is one permutation $\tau\in\mathfrak S_3$
such that, under the specified projection $p:A\to D$,
\[
                   p\alpha(\overline a_i)
                         \sim_D q(a_{\tau(i)})
                                      \quad(1\leq i\leq3).
                                                               \tag{CONE.4}
\]
In particular define the actual normal closures
\[
 \begin{aligned}
 R&=\langle\!\langle a_1,a_2,a_3\rangle\!\rangle_E,\\
 R_D&=\langle\!\langle q(a_1),q(a_2),q(a_3)\rangle\!\rangle_D.
 \end{aligned}
\]
The surjective homomorphism $h=p\alpha\eta:E\to D$, where
$\eta:E\to A$ is the source quotient, satisfies
\[
                    h(R)=R_D,
 \qquad
 \overline h:E/R=F_4\twoheadrightarrow D/R_D,
 \qquad \overline h(u_jR)=h(u_j)R_D.
                                                               \tag{CONE.5}
\]
These are maps from the stated characteristic input to the stated
lower output.  No automorphism of $D$ is asserted.

\begin{proof}
Use the actual free-product tree of the graph of groups having one
trivial central vertex, three vertices with group $C_2$, three edges
joining them to the central vertex, and four loop edges at that
central vertex.  All edge groups are trivial.  Its fundamental group
is exactly (CONE.1), and its tree action has no edge inversions.
Because the classes of the $a_i$ are nonzero in the retained
abelian quotient, $N$ meets every conjugate of a cone group trivially.
It acts freely on the tree.  Thus $N$ is the ordinary free group
of the finite graph $N\backslash T$.

The weighted vertex count of the base graph is $1+3/2$ and its
edge count is seven, so its Euler characteristic is $-9/2$.
Consequently the rank of $N$ is $r_0=1+9d/2$.
The first characteristic mod-two abelian cover has degree $2^{r_0}$,
and its free rank is exactly the displayed $r_1$, by the finite
graph-cover Euler characteristic.  The next cover has degree
$2^{r_1}$.  Therefore $[E:N_2]=m$, and the regular representation
of this actual quotient proves $\kappa_m(E)\subset N_2$.
The seven source generators also give the explicit bound
\[
                         |A|\leq(m!)^{(m!)^7}.
                                                               \tag{CONE.6}
\]

Let $Y_0=N\backslash T$, $Y_1=N_1\backslash T$, and
$Y_2=N_2\backslash T$.  They are finite connected graphs.  An
order-two element $v\in E/N_2$ acts on $Y_2$, and its projections
act on the other two graphs.  The two covering maps are equivariant
for this same abstract $C_2$ action.  On ordinary mod-two cohomology
their positive-degree maps are zero: in degree one they are the
restrictions
\[
 \operatorname{Hom}(N,\mathbb F_2)\longrightarrow
 \operatorname{Hom}(N_1,\mathbb F_2)\longrightarrow
 \operatorname{Hom}(N_2,\mathbb F_2),
                                                               \tag{CONE.7}
\]
and each target subgroup is the commutator-times-squares subgroup
of its predecessor.  Graphs have no higher cohomology.

Apply the reduced-cone Tate argument of (ICF.6) to
$D(Y_i)=\operatorname{Cone}(\mathbb F_2\to C^*(Y_i;\mathbb F_2))$.
Connectedness represents these equivariantly in degrees zero and
one.  Each covering map is zero on the ordinary-cohomology-first
$E_2$ page of the complete cyclic resolution.  Its filtration has
two pieces, so their composite is zero on Tate hypercohomology.
The tree action has no inversions; after subdivision, fixed cells
are fixed pointwise, and all other cells occur in free $C_2$ orbits.
Localization therefore identifies this Tate map with the one on
the fixed-subcomplex cones, exactly as in A.4.56.

The bottom fixed graph $Y_0^v$ is nonempty.  Otherwise all three
fixed sets would be empty by equivariance, and their reduced cones
would be the constant $\mathbb F_2$ in degree minus one with
identity transition maps, contradicting the just-proved zero map.
When the image $\overline v\in D$ is nontrivial, a fixed vertex of
$N\backslash T$ has stabilizer a conjugate of one of the actual
cone images: the other vertex stabilizers and all edge stabilizers
are trivial.  Thus $\overline v$ is conjugate to the corresponding
$q(a_j)$.  When $\overline v=1$, it belongs to every cone subgroup.
This proves the first assertion in its exact stated form.

Since $\kappa_m(E)\subset N\subset E'E^2$, the actual map gives
$H_1(A,\mathbb F_2)=H_1(E,\mathbb F_2)=\mathbb F_2^7$.
The same is true for $D$.  An automorphism $\alpha$ of $A$ therefore
sends the three named cone classes to three linearly independent
nonzero classes.  Their images in $E/N_2$ are still nontrivial
involutions.  The first assertion shows that their images in $D$
are conjugates of cone generators.  Their abelian classes are thus
three members of the list $e_1,e_2,e_3$.  Independence forces them
to be distinct, giving the one permutation $\tau$ in (CONE.4).

Finally $h$ is surjective.  The image of the normal closure $R$
is the normal closure in $D$ of these three conjugates; since their
types form the permutation $\tau$, this is exactly $R_D$.
Killing the three displayed relations in (CONE.1) leaves the four
named free generators, giving the exact source quotient $E/R=F_4$.
The universal property of that quotient supplies precisely the
surjection and generator formula in (CONE.5).
\end{proof}
